%!TEX root = TheoretischePhysikC.tex

\section{Vorlesung: Die Schrödingergleichung}



\subsection{Die Schrödingergleichung}
Von den fünf Postulaten sind die ersten vier meist schnell identifiziert; die eigentliche Arbeit steckt in Postulat 5, und analytisch lösbar ist dort nur eine Handvoll Systeme -- etwa fünf.

Gemäß Postulat 5 wird die Zeitentwicklung eines Systems durch den Zeitentwicklungsoperator
\begin{equation}
	U_t \equiv e^{-\frac{it}{\hbar}\widehat{H}}
\end{equation}
berechnet. Angenommen, dass sich das System in einem reinen Zustand $|\Psi\rangle$ mit dazugehörigem Wellenfunktion $\Psi$  befindet: Wie wird die Zeitentwicklung 
\begin{equation}
	|\Psi(t)\rangle \equiv U_t|\Psi\rangle
\end{equation}
berechnet? Im Allgemeinen wird der Zeitentwicklungsoperator $\widehat{H}$ durch einen Operator, der auf einen unendlichen dimensionalen Hilbertraum wirkt, gegeben. Solche Operatoren sind nur schwer zu diagonalisieren. Deshalb wird eine indirekte Methode verwendet, um $|\Psi(t)\rangle$ zu berechnen: Man bildet die erste Ableitung von $|\Psi(t)\rangle$ bezüglich der Zeit:
\begin{equation}
	\frac{\partial}{\partial t} |\Psi(t)\rangle = \frac{\partial}{\partial t} e^{-\frac{it}{\hbar}\widehat{H}} |\Psi\rangle = -\frac{i}{\hbar}\widehat{H} e^{-\frac{it}{\hbar}\widehat{H}} |\Psi\rangle = -\frac{i}{\hbar}\widehat{H}|\Psi(t)\rangle.
\end{equation}
Dass man so rechnen darf, liegt an der Potenzreihe: Nur die Exponentialfunktion hängt von der Zeit ab, und in ihrer Reihe verhält sich $\widehat{H}$ wie eine Zahl. Umgestellt steht dort
\begin{equation}
	i\hbar\frac{\partial}{\partial t}|\Psi(t)\rangle = \widehat{H}|\Psi(t)\rangle,
\end{equation}
die Schrödingergleichung in allgemeiner Form -- eine für jedes System. Fundamental ist bei uns aber Postulat 5; die Gleichung ist bloß das Mittel, $U_t|\Psi\rangle$ auszurechnen.

Um die Berechnungen zu vereinfachen, nehmen wir das Beispeil eines Teilchens in einer Dimension. Weil die Schrödingergleichung eine Vektorgleichung ist, dürfen wir beide Seiten mit demselben Bra multiplizieren; links vertauscht $\langle x|$ dabei mit der Zeitableitung. Wir berechnen nun das Skalarprodukt von dieser Gleichung und $|x\rangle$:
\begin{equation}
	\begin{split}
		i\hbar \langle x|\frac{\partial}{\partial t} |\Psi(t)\rangle  &= \langle x|\frac{\widehat{p}^2}{2m}|\Psi(t)\rangle \\
		&= \frac{1}{2m}\langle x|\widehat{p}^2|\Psi(t)\rangle 
	\end{split}
\end{equation}
Wir wenden jetzt die von
\begin{equation}
	\langle x|\widehat{p}^2|\Psi\rangle = -\hbar^2\frac{d^2}{dx^2}\Psi(x)
\end{equation}
gegebenen Matrixelemente von $\widehat{H}$ an und setzen die Wellenfunktion-Darstellung ein, um 
\begin{equation}\label{eq:1dfpschr}
	\eqbox{i\hbar\frac{\partial}{\partial t}\Psi(x,t) = -\frac{\hbar^2}{2m}\frac{\partial^2}{\partial x^2}\Psi(x,t)}
\end{equation}
zu erhalten. Diese Gleichung ist als die \emph{Schrödingergleichung} (für ein freies Teilchen in einer Dimension) bekannt.

Es lohnt sich, kurz bei ihrer Form zu verweilen: Streicht man das $i$, so steht dort -- bis auf ein Vorzeichen -- die Wärmeleitungsgleichung. Solche Gleichungen sind gründlich untersucht -- darin liegt der Gewinn der Herleitung: Wir dürfen den ganzen alten Werkzeugkasten anwenden.

Die Lösung der Schrödingergleichung beschreibt die Zeitentwicklung des Zustandes eines Teilchens. 
\begin{beispiel}
	Angenommen, ein Teilchen befindet sich zum Zeitpunkt $t=0$ in einem reinen Zustand -- einem \emph{Gauß'schen Wellenpaket}:
	\begin{equation}
		\Psi(x,t=0) = \sqrt{\frac{1}{\sigma\sqrt{2\pi}}}e^{-\frac14 \frac{(x-a)^2}{\sigma^2}}.
	\end{equation}
	Wie entwickelt sich die Wellenfunktion bezüglich der Zeit? Die Lösung der Schrödingergleichung (\ref{eq:1dfpschr}) in diesem Fall lautet
	\begin{equation}
		\Psi(x,t) = \sqrt{\frac{1}{\sigma(1+i(t/\tau))\sqrt{2\pi}}}e^{-\frac14 \frac{(x-a)^2}{\sigma^2(1+i(t/\tau))}},
	\end{equation}
	wobei
	\begin{equation}
		\tau \equiv \frac{2m\sigma^2}{\hbar}.
	\end{equation}	
\end{beispiel}

Wir fangen mit einer reellen Wellenfunktion an; der Faktor $1+i(t/\tau)$ macht sie komplex. Das Betragsquadrat bleibt eine Glockenkurve, nur eine immer breitere -- das Wellenpaket \emph{zerfließt}.

Bei Problemen mit Teilchen, die sich in Potentialen bewegen, ist es oft einfacher, ihre Dynamik zu berechnen, indem man, wie oben beschrieben, die Schrödingergleichung löst.


\subsection{Fourier Transformation}
Um die Dynamik freien Teilchen im Allgemeinen zu verstehen, ist es hilfreich eine neue mathematische Methode einzuführen. Rein mathematisch liefert die Fourier-Transformation eine \emph{alternative Darstellung} einer Funktion: dieselbe Information, andere Form -- zum Lösen von Gleichungen oft bequemer. In der Musik heißt die Transformierte eines Signals sein \emph{Spektrum}. Angenommen, wir haben eine Funktion $f:\mathbb{R}\rightarrow\mathbb{C}$. Die \emph{Fourier-Transformierte} der Funktion $f(x)$ wird definiert als 
\begin{equation}
	\eqbox{F(k)\equiv \frac{1}{\sqrt{2\pi}}\int_{-\infty}^\infty f(x)e^{-ikx}\, dx.}
\end{equation}
Die \emph{inverse Transformation} lautet
\begin{equation}\label{eq:invfourier}
	\eqbox{f(x)\equiv \frac{1}{\sqrt{2\pi}}\int_{-\infty}^\infty F(k)e^{ikx}\, dk.}
\end{equation}
Die Funktionen $f(x)$ und $F(k)$ sind \emph{Fourier-Paare}. Beide Funktionen beinhalten dieselbe Information. Um diese Äquivalenz zu zeigen, setzen wir die beiden Gleichungen ineinander ein:
\begin{equation}
	\begin{split}
		f(x) &= \frac{1}{\sqrt{2\pi}}\int_{-\infty}^\infty \left(\frac{1}{\sqrt{2\pi}}\int_{-\infty}^\infty f(y)e^{-iky}\, dy\right)e^{ikx}\, dk \\
		&= \int_{-\infty}^\infty f(y)\left(\frac{1}{2\pi}\int_{-\infty}^\infty e^{ik(x-y)}\, dk \right)dy \\
		&= \int_{-\infty}^\infty f(y)\delta(x-y)\,dy.
	\end{split}
\end{equation}
Hierbei haben wir in der zweiten Reihe die Darstellung der Dirac'sche Deltafunktion als 
\begin{equation}
	\delta(x) = \frac{1}{2\pi}\int_{-\infty}^\infty e^{ikx}\, dk
\end{equation}
eigesetzt. Diese Integraldarstellung der Deltafunktion ist nur \emph{eine} von mehreren möglichen; wir übernehmen sie hier ohne Beweis.

\subsection{Lösung der Schrödingergleichung für ein freies Teilchen in $1D$}
Die Schrödingergleichung eines Teilchens in einer Dimension lautet
\begin{equation}
	i\hbar\frac{\partial}{\partial t}\Psi(x,t) = -\frac{\hbar^2}{2m}\frac{\partial^2}{\partial x^2}\Psi(x,t).
\end{equation}
Wie löst man eine Gleichung dieser Form? Der erste Schritt ist bei Differentialgleichungen immer derselbe: Man probiert die Exponentialfunktion. Weil die Schrödingergleichung linear ist, suchen wir nach Lösungen der Form ebener Wellen: $\Psi(x,t) \propto e^{i(kx-\omega t)}$. Einsetzen in die Schrödingergleichung liefert
\begin{equation}
	i\hbar (-i\omega)e^{i(kx-\omega t)} = -\frac{\hbar^2}{2m}(-k^2)e^{i(kx-\omega t)}.
\end{equation}
Wir erhalten dabei eine Ebene-Welle Lösung nur dann, wenn
\begin{equation}
	\eqbox{\hbar\omega = \frac{\hbar^2k^2}{2m}.}
\end{equation} 
Diese Bedingung wird \emph{Dispersionsrelation} bezeichnet. Wenn die Dispersionsrelation erfüllt ist, dann bildet 
\begin{equation}
	\Psi_{k}(x,t) = e^{i\left(kx-\frac{\hbar k^2}{2m} t\right)}
\end{equation}
eine Familie von Lösungen, die mit $k$ bezeichnet sind. Für jedes $k$ gibt es eine solche Lösung, also unendlich viele; addiert werden sie, indem man die Summe durch ein Integral ersetzt -- dabei entdecken wir die Fourier-Transformation quasi wieder.

Wegen der Linearität der Schrödingergleichung können wir neue Lösungen nach dem Superpositionsprinzip bilden:
\begin{equation}
	\eqbox{\Psi(x,t) = \int_{-\infty}^{\infty} \Phi(k)e^{i\left(kx-\frac{\hbar k^2}{2m} t\right)}\, dk.}
\end{equation}
Um die Schrödingergleichung vollständig zu lösen, ist es nötig, die Anfangsbedingungen zu beachten:
\begin{equation}
	\Psi(x,t=0) = \int_{-\infty}^{\infty} \Phi(k)e^{ikx}\, dk.
\end{equation}
Hierbei wird $\Psi(x,0)$ vorgegeben und wir müssen für $\Phi(k)$ lösen. Um dies zu tun vergleichen wir diese Gleichung mit der inversen Transformation Gl.~(\ref{eq:invfourier}): Wir betrachten $\Phi(k)$ als die Fourier-Transformierte von $\Psi(x,0)$ und wenden die Fourier Transformation an, um
\begin{equation}
	\eqbox{\Phi(k) = \frac{1}{2\pi}\int_{-\infty}^{\infty} \Psi(x,0)e^{-ikx}\, dx}
\end{equation}
zu erhalten. Beim Vergleich stört zunächst der Vorfaktor $\frac{1}{\sqrt{2\pi}}$, der in Gl.~(\ref{eq:invfourier}) vor dem Integral steht, in der Anfangsbedingung dagegen fehlt. Die Zuordnung lautet deshalb $f(x)\equiv \Psi(x,0)$ und $F(k)\equiv \sqrt{2\pi}\,\Phi(k)$, und die Hintransformation liefert
\begin{equation}
	\sqrt{2\pi}\,\Phi(k) = \frac{1}{\sqrt{2\pi}}\int_{-\infty}^{\infty} \Psi(x,0)e^{-ikx}\, dx.
\end{equation}
Der Vorfaktor $\frac{1}{2\pi}$ oben ist also nicht geraten, sondern Buchhaltung. Für die Fourier-Transformation sind viele Konventionen im Umlauf; prüfen Sie in fremden Texten zuerst nach, welche gilt.

Warnung: Eine ebene Welle ist \emph{nicht} ein Element von $L^2(\mathbb{R})$, weil die Wellenfunktion $\Psi_{k}(x,t)$ nicht quadratintegriebar ist. Die Funktion $\Psi_{k}(x,t)$ bildet zwar eine Lösung, ist aber selbst nicht ein Element von dem Hilbertraum $\mathcal{H}$. Diese Beobachtung schließt nicht die Möglichkeit aus, dass wir durch lineare Kombinationen von ebenen Wellen quadratintegriebare Funktionen bilden können. Man sieht das sofort an der Wahrscheinlichkeitsdichte: $|\Psi_{k}(x,t)|^2 = 1$, überall gleich groß, und das Integral darüber divergiert -- auch keine Normierungskonstante macht es endlich.

Obwohl die ,,stationäre`` ebene Welle $\Psi_{k}(x) \equiv e^{ikx}$ keinem Vektor im Hilbertraum entspricht,  können wir dennoch eine Ket-Notation einführen: Der ,,Ket-Vektor`` $|\Psi_{k}\rangle$ steht für die ebene Welle. Die dazugehörige Wellenfunktion wird durch
\begin{equation}
	\langle x| \Psi_{k}\rangle = e^{ikx}
\end{equation} 
gegeben.


\subsection{Impuls-Raum}
Wir haben bisher nur Messungen von der Teilchenposition berücksichtigt. Wie kann aus der Wellenfunktion $\Psi(x,t)$ die Statistik für \emph{Impulsmessungen} oder für Größen, die nur vom Impuls abhängen, gewonnen werden? Für eine stationäre ebene Welle $\Psi(x,t) = e^{ikx}$ ist der Impuls durch
\begin{equation}
	\langle x|\widehat{p}|\Psi_k\rangle \overset{!}{=}  -i\hbar \frac{d}{dx} e^{ikx}  \overset{!}{=} \hbar k   e^{ikx} \overset{!}{=} \hbar k \langle x|\Psi_k\rangle
\end{equation}
zu interpretieren: Obwohl der Ket $|\Psi_k\rangle$ nicht als Element in $\mathcal{H}$ existiert, behandeln wir $|\Psi_k\rangle$ trotzdem als einen fiktiven Vektoren. Nach dieser Berechnung ist es sinnvoll $|\Psi_k\rangle$ als eine Art ,,Eigenvektor`` für $\widehat{p}$ zu interpretieren. Daraus ergibt sich, dass eine ebene Welle $\Psi(x,t) = e^{ikx}$ einen bestimmen Impuls von $\hbar k$ hat. Der Vektorgröße $|\Psi_k\rangle$ spielt für $\widehat{p}$ genau diesselbe Rolle, dass $|x\rangle$ für $\widehat{x}$ spielt.

\begin{intuition}
Denken Sie an einen gewöhnlichen Vektor, zerlegt nach zwei verschiedenen Basen: derselbe Vektor, zwei völlig verschiedene Zahlenlisten. Genauso hier -- die Koeffizienten von $|\Psi\rangle$ heißen in der Positionsbasis \emph{Wellenfunktion}, in der Basis der ebenen Wellen \emph{Impulswellenfunktion}.
\end{intuition}

Mit Hilfe dieser Überlegungen führen wir jetzt die folgende Wellenfunktion im Impulsraum ein:
\begin{equation}
	\Phi(p,t) \equiv \frac{1}{\sqrt{2\pi\hbar}} \int_{-\infty}^\infty \Psi(x,t)e^{-i\frac{px}{\hbar}}\,dx.	
\end{equation} 
Man kann sehen, dass diese neue \emph{Impulswellenfunktion} einfach durch die Fourier-Transformation der ursprünglichen Ortswellenfunktion gegeben wird. Daraus schließen wir, dass die Impulswellenfunktion die gleiche Information enthält wie die Ortswellenfunktion; sie beschreiben in unterschiedlichen Darstellungen den gleichen Zustand des Teilchens. Angenommen, wir haben nur die Impulswellenfunktion, können wir die Ortswellenfunktion  mit der inversen Fourier transformation rekonstruieren:
\begin{equation}
	\Psi(x,t) \equiv \frac{1}{\sqrt{2\pi\hbar}} \int_{-\infty}^\infty \Phi(p,t)e^{i\frac{px}{\hbar}}\,dp.	
\end{equation} 
Der Vorfaktor $\frac{1}{\sqrt{2\pi\hbar}}$ und das $\hbar$ im Exponenten wirken unmotiviert; sie sind der Preis für eine schöne Analogie: Mit dieser Konvention steht der Impulsmittelwert nachher genau dort, wo im Ortsraum $\langle\widehat{x}\rangle = \int x|\Psi(x,t)|^2\,dx$ steht.

% fig_vl08_fourier.tex — Das Fourier-Paar: dasselbe Wellenpaket im Orts-
% und im Impulsraum, und die reziproke Breite der beiden Darstellungen.
%
% Follows the conventions of fig_vl11_kasten.tex.
%
%   * geometry from the closed form, not freehand.  The lecture's packet
%         Psi(x,0) = (sigma sqrt(2 pi))^{-1/2} exp(-(x-a)^2 / (4 sigma^2))
%     has the momentum wave function (Fourier transform, p = hbar k)
%         Phi(p)  proportional to exp(-sigma^2 p^2 / hbar^2),
%     so the two densities are Gaussians of standard deviation
%         sigma_x = sigma        and      sigma_p = hbar / (2 sigma).
%     Every curve, every bar length and every peak height in the plate is
%     evaluated from those two numbers:
%         |Psi|^2 peak = 1/(sigma_x sqrt(2 pi)),
%         |Phi|^2 peak = 1/(sigma_p sqrt(2 pi)),
%     which is why the tallest curve of (a) is the flattest of (b).  The
%     three states are sigma_x = sigma/2, sigma, 2 sigma, spaced by a
%     factor of two so that the reciprocal factor of two in (b) can be
%     read off the ruler rather than believed.
%
%   * plotted are the *densities* |Psi|^2 and |Phi|^2, because those are
%     what the lecture integrates in <p> = int p |Phi(p,t)|^2 dp;
%   * colour encodes the state (which packet), nothing else — the same
%     three colours in the same order appear in both panels, so the eye
%     can carry a packet across the Fourier transform;
%   * the +/- sigma rulers sit in a strip above the curves, clear of the
%     data: their lengths are the whole point of the plate, and their
%     order reverses from (a) to (b);
%   * dotted grey drop lines take the middle packet's ruler down onto its
%     own flank in both panels, so the bars are visibly *measured* off a
%     curve rather than asserted; that packet is the worked example, the
%     other two are the trend;
%   * the grey double arrow in the gap carries the only thing the two
%     panels do not say separately: they are one state, not two;
%   * axis units are sigma and hbar/sigma, so the plate carries no
%     arbitrary numbers, and the y axes are unnumbered — the message is
%     the width, not the height.
%
% Panel (a) three packets in position space, narrow to wide.
% Panel (b) the same three in momentum space, wide to narrow, with the
%           product sigma_x sigma_p = hbar/2 that the two panels measure.

\begin{center}
\begin{tikzpicture}

% Dimension ruler: bar–arrow–arrow–bar, the width of one packet.
\tikzset{
  dimbar/.style={line width=0.7pt, >={Stealth[length=4pt,width=3pt]}, |<->|},
}

% ─────────────────────────────────────────────────────────────
%  (a)  Position space
% ─────────────────────────────────────────────────────────────
\begin{axis}[
    name=posspace,
    tpcplot,
    width=6.4cm, height=5.9cm,
    grid=none,
    xmin=-5, xmax=5, ymin=0, ymax=2.62,
    xtick={-4,-2,0,2,4}, ytick=\empty,
    xlabel={$x/\sigma$},
    ylabel={$|\Psi(x)|^2$},
    every axis x label/.style={at={(axis description cs:0.5,-0.13)},
                               anchor=north, font=\small},
    every axis y label/.style={at={(axis description cs:-0.045,0.5)},
                               rotate=90, anchor=south, font=\small},
    title={(a) \tpcifenglish{position space}{Ortsraum}},
    domain=-5:5, samples=201,
    clip=false,
]
  % |Psi|^2 = exp(-x^2/(2 sigma_x^2)) / (sigma_x sqrt(2 pi))
  \addplot[wavefn, spacecadet]   {1.59577*exp(-2*x^2)};       % sigma_x = sigma/2
  \addplot[wavefn, cgblue]       {0.79788*exp(-0.5*x^2)};     % sigma_x = sigma
  \addplot[wavefn, birrenorange] {0.39894*exp(-0.125*x^2)};   % sigma_x = 2 sigma

  % Width rulers, +/- sigma_x, in the state's own colour
  \draw[dimbar, spacecadet]   (axis cs:-0.5,2.40) -- (axis cs:0.5,2.40);
  \draw[dimbar, cgblue]       (axis cs:-1,  2.15) -- (axis cs:1,  2.15);
  \draw[dimbar, birrenorange] (axis cs:-2,  1.90) -- (axis cs:2,  1.90);

  \node[lbl, font=\scriptsize, text=spacecadet,   anchor=west]
    at (axis cs:0.65,2.40) {$\sigma/2$};
  \node[lbl, font=\scriptsize, text=cgblue,       anchor=west]
    at (axis cs:1.15,2.15) {$\sigma$};
  \node[lbl, font=\scriptsize, text=birrenorange, anchor=west]
    at (axis cs:2.15,1.90) {$2\sigma$};

  % The ruler really is read off the curve: drop the middle packet's
  % +/- sigma_x onto its own flank, where |Psi|^2 has fallen to e^{-1/2}
  % of its peak.  The same packet is worked in panel (b).
  \draw[classicallim] (axis cs:-1,2.15) -- (axis cs:-1,0.484);
  \draw[classicallim] (axis cs: 1,2.15) -- (axis cs: 1,0.484);

  \node[lbl, font=\scriptsize, text=birrengray, anchor=west]
    at (axis cs:-4.9,2.54) {$\pm\sigma_x$};
\end{axis}

% ─────────────────────────────────────────────────────────────
%  (b)  Momentum space — the same three packets
% ─────────────────────────────────────────────────────────────
\begin{axis}[
    name=momspace,
    tpcplot,
    at={(posspace.south east)}, anchor=south west, xshift=19mm,
    width=6.4cm, height=5.9cm,
    grid=none,
    xmin=-2.5, xmax=2.5, ymin=0, ymax=2.62,
    xtick={-2,-1,0,1,2}, ytick=\empty,
    xlabel={$p\,\sigma/\hbar$},
    ylabel={$|\Phi(p)|^2$},
    every axis x label/.style={at={(axis description cs:0.5,-0.13)},
                               anchor=north, font=\small},
    every axis y label/.style={at={(axis description cs:-0.045,0.5)},
                               rotate=90, anchor=south, font=\small},
    title={(b) \tpcifenglish{momentum space}{Impulsraum}},
    domain=-2.5:2.5, samples=201,
    clip=false,
]
  % |Phi|^2 with sigma_p = hbar/(2 sigma_x): the order of the widths flips
  \addplot[wavefn, spacecadet]   {0.39894*exp(-0.5*x^2)};     % sigma_p = hbar/sigma
  \addplot[wavefn, cgblue]       {0.79788*exp(-2*x^2)};       % sigma_p = hbar/(2 sigma)
  \addplot[wavefn, birrenorange] {1.59577*exp(-8*x^2)};       % sigma_p = hbar/(4 sigma)

  \draw[dimbar, spacecadet]   (axis cs:-1,   2.40) -- (axis cs:1,   2.40);
  \draw[dimbar, cgblue]       (axis cs:-0.5, 2.15) -- (axis cs:0.5, 2.15);
  \draw[dimbar, birrenorange] (axis cs:-0.25,1.90) -- (axis cs:0.25,1.90);

  \node[lbl, font=\scriptsize, text=spacecadet,   anchor=west]
    at (axis cs:1.15,2.40) {$\hbar/\sigma$};
  \node[lbl, font=\scriptsize, text=cgblue,       anchor=west]
    at (axis cs:0.65,2.15) {$\hbar/(2\sigma)$};
  \node[lbl, font=\scriptsize, text=birrenorange, anchor=west]
    at (axis cs:0.40,1.90) {$\hbar/(4\sigma)$};

  \draw[classicallim] (axis cs:-0.5,2.15) -- (axis cs:-0.5,0.484);
  \draw[classicallim] (axis cs: 0.5,2.15) -- (axis cs: 0.5,0.484);

  \node[lbl, font=\scriptsize, text=birrengray, anchor=west]
    at (axis cs:-2.45,2.54) {$\pm\sigma_p$};

  % What the two rulers measure together
  \node[lbl, font=\small, text=cgblue, anchor=west]
    at (axis cs:-2.45,1.30) {$\sigma_x\,\sigma_p = \hbar/2$};
\end{axis}

% The two panels are one state, not two: the transform lives in the gap.
\coordinate (fmid) at ($(posspace.north east)!0.5!(momspace.north west)+(0,-11mm)$);
\draw[{Stealth[length=4pt,width=3pt]}-{Stealth[length=4pt,width=3pt]},
      birrengray, line width=0.7pt]
  ($(fmid)+(-4mm,0)$) -- ($(fmid)+(4mm,0)$);
\node[lbl, font=\scriptsize, text=birrengray, anchor=south]
  at ($(fmid)+(0,0.7mm)$) {$\mathcal{F}$};

\end{tikzpicture}
\end{center}


Vergleichen Sie in der Abbildung die drei Balken über den Kurven: Halbiert man die Breite im Ortsraum, so verdoppelt sie sich im Impulsraum. Gleiche Farbe heißt dasselbe Paket -- ihre Reihenfolge kippt, das Produkt $\sigma_x\sigma_p$ nicht.


\subsubsection{Statistik von Impulsmessungen}
Die Statistik der Impulsmessungen kann von der Positionsraum-Wellenfunktion abgelesen werden:
\begin{align}
	\langle \widehat{p}\rangle &= \int_{-\infty}^{\infty} \Psi^*(x,t) \left(-i\hbar \frac{\partial}{\partial x} \Psi(x,t)\right)\,dx \\
	\langle \widehat{p}^2\rangle &= \int_{-\infty}^{\infty} \Psi^*(x,t) \left(-\hbar^2 \frac{\partial^2}{\partial x^2} \Psi(x,t)\right)\,dx.
\end{align}
Alternativ können wir die Impulsstatistik mit Hilfe der Impulswellenfunktion berechnen, z.B.:
\begin{equation}
	\begin{split}
		\langle \widehat{p}\rangle &= \frac{1}{2\pi\hbar}\int_{-\infty}^{\infty}\int_{-\infty}^{\infty}\int_{-\infty}^{\infty} \Phi(p',t)e^{-i\frac{p'x}{\hbar}} \left(-i\hbar \frac{\partial}{\partial x} \Phi^*(p,t)e^{i\frac{px}{\hbar}}\right)\,dxdpdp' \\
		&= \frac{1}{2\pi\hbar}\int_{-\infty}^{\infty}\int_{-\infty}^{\infty}\int_{-\infty}^{\infty} \Phi^*(p',t)e^{-i\frac{p'x}{\hbar}} \left(p \Phi(p,t)e^{i\frac{px}{\hbar}}\right)\,dxdpdp' \\
		&= \int_{-\infty}^{\infty}\int_{-\infty}^{\infty} p \Phi^*(p',t) \Phi(p,t) \left(\frac{1}{2\pi\hbar}\int_{-\infty}^{\infty} e^{i\frac{(p-p')x}{\hbar}}\,dx\right)\,dpdp' \\
		&= \int_{-\infty}^{\infty}\int_{-\infty}^{\infty} p \Phi^*(p',t) \Phi(p,t) \left(\frac{1}{2\pi}\int_{-\infty}^{\infty} e^{i{(p-p')x}}\,dx'\right)\,dpdp' \\
		&= \int_{-\infty}^{\infty}\int_{-\infty}^{\infty} p \Phi^*(p',t) \Phi(p,t) \delta(p-p') \,dpdp'\\
		&= \int_{-\infty}^{\infty} p |\Phi(p,t)|^2 \,dp.
	\end{split}
\end{equation}
Nach einem ähnlichen Verfahren erhalten wir:
\begin{equation}
		\langle \widehat{p}^2\rangle = \int_{-\infty}^{\infty} p^2 |\Phi(p,t)|^2 \,dp.
\end{equation}
Das heißt, die Statistik der Impulsmessungen ist wesentlich einfacher bezüglich Impulswellenfunktionen zu berechnen. Vergleichen Sie den Aufwand: erst ableiten, mit $\Psi^*$ multiplizieren, über $x$ integrieren -- oder nur $|\Phi(p,t)|^2$ mit dem Gewicht $p$ integrieren.


\subsection{Ein Teilchen in drei Dimensionen}
Eigentlich ist das Problem schon gelöst: Ein freies Teilchen in drei Dimensionen ist nichts anderes als drei Kopien eines Teilchens in einer Dimension und über den Separationsansatz mitgelöst. Ein Teilchen in drei Dimensionen wird in der Quantenmechanik wie folgt beschrieben
\begin{enumerate}
	\item Hilbertraum: $L^2(\mathbb{R}^3)$.
	\item Präparationen von reinen Zuständen $\rho \equiv |\Psi\rangle\langle\Psi|$: Wellenfunktionen $\Psi:\mathbb{R}^3\rightarrow \mathbb{C}$, wobei
	\begin{equation}
		|\Psi\rangle \equiv \iiint_{-\infty}^\infty \Psi(x,y,z) |x,y,z\rangle \, dxdydz
	\end{equation}
	und
	\begin{equation}
		\langle x',y',z'|x,y,z\rangle = \delta(x'-x)\delta(y'-y)\delta(z'-z).
	\end{equation}
	Damit eine solche Funktion als Wellenfunktion durchgeht, muss sie -- genau wie in einer Dimension -- auf eins normiert sein:
	\begin{equation}
		\langle\Psi|\Psi\rangle = \iiint_{-\infty}^\infty |\Psi(x,y,z)|^2 \, dxdydz = 1.
	\end{equation}
	\item Detektoren: Sei $S\subset \mathbb{R}^3$ ein Bereich. Der Messoperator, der die Messung, ob sich das Teilchen in einer Region S befindet, modelliert, ist durch 
	\begin{equation}
		F_{S} \equiv \iiint_S  	|x,y,z\rangle \langle x,y,z|\, dxdydz
	\end{equation} 
	gegeben.
	\item Ansprechhäufigkeit des Detektors $F_S$: 
	\begin{equation}
		\mathbb{P}[\text{,,$(x,y,z)\in S$``}] \equiv \tr(\rho F_S) = \iiint_S |\Psi(x,y,z)|^2 \, dxdydz.
	\end{equation}
	\item Hamiltonoperator:
	\begin{equation}
		\begin{split}
			\widehat{H} \equiv \frac{\widehat{P}^2_{\text{ges.}}}{2m} &= \frac{\widehat{p}_x^2}{2m}+\frac{\widehat{p}_y^2}{2m}+\frac{\widehat{p}_z^2}{2m} \\ &= -\frac{\hbar^2}{2m}\frac{\partial^2}{\partial x^2}-\frac{\hbar^2}{2m}\frac{\partial^2}{\partial y^2}-\frac{\hbar^2}{2m}\frac{\partial^2}{\partial z^2}.
		\end{split}
	\end{equation}
\end{enumerate}


\subsection{Die Quantisierungsregeln für das Teilchen}
Die Ableitung des Hamiltonoperators für das freie Teilchen ist ein Beispiel für ein Verfahren, wie
aus klassischen Größen ,,entsprechende`` quantenmechanische gewonnen werden können: Dabei werden aus Funktionen der klassischen Variablen Ort und Impuls Operatoren im Raum der Wellenfunktionen. Die folgende Korrespondenz-Tabelle stellt ein Paar Quantisierungsregeln für das Teilchen (in drei Dimensionen) dar
\begin{center}
\begin{tabular}{lcc}    \toprule
{\bf klassisch} & & {\bf quantenmechanischer Operator wirkt nach}  \\\midrule
Ort ($j$-Koord.): &$x_j$ & $\widehat{x}_j \Psi(x,y,z) = x_j \Psi(x,y,z)$  \\ 
\rowcolor{isabelline} Impuls ($j$-Koord.):& $p_j$ & $\widehat{p}_j \Psi(x,y,z) = -i \hbar \frac{\partial}{\partial x_j}\Psi(x,y,z)$  \\ 
kinetische Energie: & $K$ & $\widehat{K}\Psi(x,y,z) = -\frac{\hbar^2}{2m} \nabla^2 \Psi(x,y,z)$ \\
\rowcolor{isabelline} potentielle Energie & $V(x,y,z)$ & $\widehat{V} \Psi(x,y,z) = V(x,y,z)\Psi(x,y,z)$ \\
gesamte Energie & $E$ & $\widehat{H}\Psi(x,y,z) = \left(-\frac{\hbar^2}{2m} \nabla^2+ V(x,y,z)\right)\Psi(x,y,z)$ \\
\rowcolor{isabelline} Drehimpuls & $\mathbf{L}$ & $\widehat{L}_1\Psi(x,y,z) = -i \hbar (y \frac{\partial}{\partial z} - z \frac{\partial}{\partial y})\Psi(x,y,z)$  \\
allgemein & $f$ & ??? \\
\bottomrule
\end{tabular}
\end{center}

Kurz gesagt also, lautet die Regel: Ersetze in jedem klassischen Ausdruck die Ortskoordinaten durch Multiplikationsoperatoren, und die Impulskoordinaten durch ($-i\hbar \times$) Differentiationsoperatoren. Diese Regel ist recht nützlich, und man sollte sie sich unbedingt merken. Davor steht immer ein physikalischer Schritt: Zuerst muss man die \emph{klassische} Energie verstehen, erst daraus gewinnt man den Operator. Und die Regel ist kein wohldefiniertes Verfahren, sondern eine Faustregel -- jedes Mal eine Vermutung, über die erst der Vergleich mit dem Experiment entscheidet. Es gibt aber eine entscheidende Einschränkung: Klassisch ist der Ausdruck 
\begin{equation*}
	p_j x_k - x_k p_j = 0.
\end{equation*}
Mit den oben eingeführten Operatoren dagegen erhält man
\begin{equation}
	\begin{split}
		-i \hbar \left(\frac{\partial}{\partial x_j}x_k - x_k \frac{\partial}{\partial x_j}\right)\Psi(x,y,z) &= -i \hbar \left(x_k\frac{\partial}{\partial x_j}  + \delta_{jk} - x_k \frac{\partial}{\partial x_j} \right)\Psi(x,y,z) \\
		&= -i\hbar \delta_{jk} \Psi(x,y,z).
	\end{split}
\end{equation}
Das heißt, der entsprechende quantenmechanische Operator wirkt wie $-i\hbar\delta_{jk}\mathbb{I}$. Zusammengefasst schreiben wir
\begin{equation}
	[\widehat{p}_j,\widehat{x}_k]\equiv\widehat{p}_j\widehat{x}_k - \widehat{x}_k  \widehat{p}_j = -i\hbar\delta_{jk}\mathbb{I},
\end{equation} 
wobei $[\widehat{A},\widehat{B}] = \widehat{A}\widehat{B} - \widehat{B}\widehat{A}$ den \emph{Kommutator} zweier Operatoren bezeichnet. Diese Operatorgleichung wird als kanonische Vertauschungsrelationen bezeichnet. Der ganze Unterschied ist dabei die Leibniz-Regel: Einmal multipliziert man zuerst mit $x_k$ und leitet dann das Produkt ab, einmal leitet man zuerst ab und multipliziert danach -- der so entstehende Term $\delta_{jk}$ ist, bis auf den Faktor $-i\hbar$, der ganze Kommutator.

Für das Ausfüllen der obigen Korrespondenztabelle bedeuten sie einfach, dass es keine allgemein gültige ,,Quantisierungsvorschrift`` für klassische Größen gibt. Die konkreten Einträge in der Tabelle sind übrigens gerade solche Funktionen von $x$ und $p$, bei denen solche \emph{Operatorordnungsprobleme} keine Rolle spielen. Quantenmechanisch stehen für die klassische Größe $x\,p_x$ dagegen drei Kandidaten nebeneinander -- $\widehat{x}\widehat{p}_x$, $\widehat{p}_x\widehat{x}$ und das symmetrisierte Mittel --, und nichts zeichnet einen davon aus. Die Differenz $\widehat{x}\widehat{p}_x-\widehat{p}_x\widehat{x}$ ist klassisch exakt null, quantenmechanisch mit $\hbar\approx 10^{-34}\,\mathrm{J\,s}$ nur winzig -- aber eben nicht null, und dieser Unterschied ist beobachtbar. Er läuft unter dem Stichwort \emph{Unschärfe}.

